\documentclass[aspectratio=169,10pt]{beamer}
\usetheme{Madrid}
\usecolortheme{default}
\usepackage{booktabs,graphicx,amsmath,array}
\definecolor{passgreen}{RGB}{0,110,60}
\definecolor{failred}{RGB}{180,30,30}
\setbeamertemplate{navigation symbols}{}
\setbeamertemplate{footline}[frame number]
\newcommand{\evidence}[2]{\IfFileExists{assets/#1}{\includegraphics[width=\linewidth,height=4.2cm,keepaspectratio]{assets/#1}}{\fbox{\parbox[c][3.5cm][c]{0.93\linewidth}{\centering #2\\[1em]\scriptsize Add image: \texttt{\detokenize{assets/#1}}}}}}
\title{Dense-to-Sparse Crowd Localization}
\subtitle{Progress: CPU Reference, FPGA Implementation, and Validation}
\author{Ransara}
\date{20 October 2026}
\begin{document}
\begin{frame}\titlepage\end{frame}

\begin{frame}{Objective and current scope}
\begin{itemize}
\item Convert confidence map $P(x,y)$ and threshold map $T(x,y)$ into sparse head coordinates.
\item Compare an OpenCV CPU baseline with a margin-gated, streaming FPGA decoder.
\item Accelerate the \textbf{decoder}; neural-network inference is outside this comparison.
\item Current board build: DE2-115, $32\times32$ maximum input, 64 label slots, USB-Blaster communication.
\item Current evidence: CPU reference checks pass; physical FPGA results contain mismatches and timing violations.
\end{itemize}
\begin{block}{Data limitation}The reported 20-photo experiment uses annotation-driven synthetic maps, one crop per distinct photo. It does not measure trained-model localization accuracy.\end{block}
\end{frame}

\begin{frame}{Baseline and proposed decoder}
\begin{columns}[T]
\column{0.48\textwidth}
\begin{block}{Baseline: CPU / OpenCV}
\begin{enumerate}
\item Direct comparison $B=\mathbf{1}[P\geq T]$.
\item OpenCV connected-component labeling.
\item Filter small components; compute centroids.
\end{enumerate}
\end{block}
\column{0.48\textwidth}
\begin{block}{Proposed: margin-gated streaming CCL}
\begin{enumerate}
\item $P<0.2$: background; $P\geq0.7$: foreground.
\item Compare against $T$ only in the ambiguous band.
\item Streaming labels, equivalence resolution, component accumulation and closing.
\end{enumerate}
\end{block}
\end{columns}
\medskip
Both use 4-connectivity and minimum area 3. Identical input quantization is required for the board comparison.
\[
(x_c,y_c)=\left(\frac{\sum x}{\mathrm{area}},\frac{\sum y}{\mathrm{area}}\right)
\]
\end{frame}

\begin{frame}{Code upgrades: one-slide overview}
\small
\begin{tabular}{p{0.37\linewidth}p{0.56\linewidth}}
\toprule
Change & Purpose \\
\midrule
Shared packet engine and JTAG UART adapter & Transfer maps and results through USB-Blaster; retain the decoder. \\
JTAG Quartus project and IP-generation Tcl & Generate Intel JTAG UART IP and integrate the board wrapper. \\
Python JTAG Atlantic wrapper & Resolve Linux dependencies and C++ exports; transfer binary packets. \\
\texttt{validate\_fpga.py} & Compare exact CPU/FPGA components, print head coordinates, save timings and JSON. \\
Sample and tile preparation scripts & Prepare more photos or cover one full map with independent tiles. \\
Small board configuration & Reduce input limits to $32\times32$ and label slots to 64 for bring-up. \\
\bottomrule
\end{tabular}
\end{frame}

\begin{frame}{CPU-only validation: what validate.py establishes}
\begin{itemize}
\item Margin-gated binarization equals direct $P\geq T$ in 20 random trials.
\item Python streaming CCL is compared with OpenCV on the same saved binary masks.
\item This stage runs \textbf{both implementations on the CPU}; no FPGA is involved.
\end{itemize}
\begin{center}
\begin{tabular}{lrrc}
\toprule
Sample & Python components & OpenCV components & Centroid agreement \\
\midrule
0002 & 29 & 29 & PASS \\
0004 & 87 & 87 & PASS \\
0005 & 5 & 5 & PASS \\
\bottomrule
\end{tabular}
\end{center}
\small Rechecked for this presentation. The Python reference is a correctness model, not an optimized CPU implementation. These full-map tests differ from the later 32$\times$32 board tests.
\end{frame}

\begin{frame}{Physical FPGA implementation and validation path}
\begin{center}
\fbox{Python host}\ $\longrightarrow$\ \fbox{USB-Blaster / JTAG UART}\par\medskip
$\longrightarrow$\ \fbox{Avalon adapter / packet engine}\par\medskip
$\longrightarrow$\ \fbox{Margin-gated CCL}\ $\longrightarrow$\ \fbox{Component records}
\end{center}
\begin{itemize}
\item Quartus: synthesis $\rightarrow$ fitting $\rightarrow$ assembly $\rightarrow$ timing analysis.
\item Program the generated \texttt{.sof}; send identical 8-bit $P/T$ inputs to the CPU and board.
\item Compare sorted $(\mathrm{area},\sum x,\sum y)$ records, not just head counts.
\item A successful configuration download does \textbf{not} establish timing closure or algorithm correctness.
\end{itemize}
\end{frame}

\begin{frame}{Compilation problems and actions taken}
\small
\begin{tabular}{p{0.29\linewidth}p{0.31\linewidth}p{0.32\linewidth}}
\toprule
Problem & Diagnosis & Action / status \\
\midrule
Undefined \texttt{jtag\_uart\_system} & Missing generated synthesis IP / QIP & Generate Qsys HDL and include its QIP. \\
Qsys commands unavailable & Installation tools absent from PATH & Select a consistent Quartus environment. \\
Tcl / device errors & Unsupported \texttt{file normalize}; device suffix & Remove unsupported path operation; use EP4CE115F29C7. \\
Synthesis exceeded 5 hours & Large label storage and dependent combinational lookups & Reduce 2048 labels to 64 and frame limit to $32\times32$. \\
Compilation completes, timing fails & Setup slack $-41.265$ ns at 50 MHz & \textbf{Unresolved:} pipeline / sequence lookups or use a slower physical clock. \\
\bottomrule
\end{tabular}
\end{frame}

\begin{frame}{Compilation evidence: insert screenshots}
\begin{columns}[T]
\column{0.49\textwidth}\centering
\evidence{long_compilation.png}{Original build: synthesis exceeded five hours}
\column{0.49\textwidth}\centering
\evidence{timing_report.png}{Small build: final timing report showing negative setup slack}
\end{columns}
\medskip\small The smaller build improved compilation feasibility; it did not resolve the 50 MHz timing requirement.
\end{frame}

\begin{frame}{Host integration problems and first board test}
\begin{itemize}
\item Missing shared-library symbols: preload matching Quartus clients in the required order.
\item Missing API names: bind Quartus 25.1 C++ exports in the Python wrapper.
\item JTAG UART node detected and connection established.
\item Deterministic test passed: two area-4 components at $(0.5,0.5)$ and $(5.5,0.5)$; 390 active cycles.
\item Realistic crops exposed discrepancies that this simple test did not reveal.
\end{itemize}
\begin{block}{Lesson}A simple smoke test establishes basic dataflow. Diverse masks and timing closure are necessary to establish reliable operation.\end{block}
\end{frame}

\begin{frame}{20-photo experiment: protocol and parameter definitions}
\small
\begin{itemize}
\item 20 distinct JHU validation photos; one annotation-centered $32\times32$ crop per photo.
\item Synthetic $P/T$ maps; shared unsigned round-to-nearest 8-bit quantization.
\item CPU: direct thresholding + OpenCV, 4-connectivity, minimum area 3.
\item FPGA: margin-gated thresholding + streaming CCL with matching settings.
\end{itemize}
\begin{tabular}{p{0.27\linewidth}p{0.66\linewidth}}
\toprule
Parameter & Interpretation \\
\midrule
CPU / FPGA heads & Number of components surviving the area filter. \\
CPU time & Median of repeated CPU runs: quantization, binarization, OpenCV CCL. \\
FPGA round-trip & Host quantization, transfer and result reception; excludes opening JTAG. \\
Active cycles & Hardware counter excluding communication stalls as defined by the RTL. \\
Estimated core time & Cycles / assumed 50 MHz; reliable operation requires timing closure. \\
PASS / FAIL & Exact agreement / disagreement of component areas and coordinate sums. \\
\bottomrule
\end{tabular}
\end{frame}

\begin{frame}{20-photo results: exact component agreement}
\centering\scriptsize
\setlength{\tabcolsep}{9pt}
\renewcommand{\arraystretch}{0.95}
\begin{tabular}{lrrrrc}
\toprule
Photo & CPU heads & FPGA heads & CPU ($\mu$s) & Round-trip (s) & Result \\
\midrule
0003 & 0 & 0 & 15.9 & 1.886 & \textcolor{passgreen}{PASS} \\
0007 & 1 & 1 & 16.0 & 1.554 & \textcolor{passgreen}{PASS} \\
0009 & 1 & 1 & 53.5 & 1.361 & \textcolor{passgreen}{PASS} \\
0011 & 1 & 1 & 49.8 & 0.955 & \textcolor{passgreen}{PASS} \\
0013 & 0 & 0 & 46.4 & 1.466 & \textcolor{passgreen}{PASS} \\
0042 & 1 & 1 & 41.1 & 1.154 & \textcolor{passgreen}{PASS} \\
0056 & 0 & 0 & 50.0 & 1.056 & \textcolor{passgreen}{PASS} \\
0058 & 1 & 1 & 50.0 & 1.055 & \textcolor{passgreen}{PASS} \\
0068 & 4 & 2 & 101.4 & 1.759 & \textcolor{failred}{FAIL} \\
0070 & 1 & 1 & 49.9 & 1.560 & \textcolor{passgreen}{PASS} \\
0085 & 1 & 1 & 51.5 & 1.770 & \textcolor{passgreen}{PASS} \\
0092 & 0 & 0 & 49.7 & 1.058 & \textcolor{passgreen}{PASS} \\
0105 & 1 & 1 & 41.7 & 1.056 & \textcolor{passgreen}{PASS} \\
0111 & 0 & 0 & 93.6 & 1.055 & \textcolor{passgreen}{PASS} \\
0116 & 4 & 4 & 44.7 & 0.957 & \textcolor{failred}{FAIL} \\
0125 & 1 & 1 & 49.8 & 1.466 & \textcolor{passgreen}{PASS} \\
0133 & 1 & 1 & 49.9 & 1.069 & \textcolor{passgreen}{PASS} \\
0134 & 3 & 2 & 43.7 & 0.965 & \textcolor{failred}{FAIL} \\
0136 & 2 & 2 & 106.1 & 0.953 & \textcolor{passgreen}{PASS} \\
0138 & 1 & 1 & 48.5 & 0.845 & \textcolor{passgreen}{PASS} \\
\bottomrule
\end{tabular}
\par\medskip 17/20 exact matches; 18/20 head-count matches. These are decoder agreement rates, not detection accuracy.
\end{frame}

\begin{frame}{Failure analysis: equal counts are insufficient}
\small
\begin{tabular}{p{0.10\linewidth}p{0.13\linewidth}p{0.66\linewidth}}
\toprule
Photo & CPU/FPGA & Observed discrepancy \\
\midrule
0068 & 4 / 2 & One record matches; three CPU components are represented by one FPGA component with a different total area. \\
0116 & 4 / 4 & Identical centroids; the component at crop $(16,16)$ has CPU area 5 versus FPGA area 3. \\
0134 & 3 / 2 & Two CPU components appear as one FPGA component; their combined area and sums match the merged record. \\
\bottomrule
\end{tabular}
\medskip
\begin{alertblock}{Interpretation}
Timing violations are a plausible contributor, not a proven sole cause. Reproduce failures in simulation and at a timing-clean physical clock to distinguish logic defects from timing failures.
\end{alertblock}
\end{frame}

\begin{frame}{Head locations and demonstration video}
\begin{columns}[T]
\column{0.53\textwidth}
\begin{itemize}
\item Crop coordinates: position within the tile.
\item Source coordinates: crop position plus tile origin.
\item Example: crop $(16,16)$ + origin $(961,561)$ = source $(977,577)$.
\item These are decoded centroids, not automatically ground-truth annotations.
\end{itemize}
\medskip
\href{run:assets/demo.mp4}{\beamergotobutton{Open 20-photo demonstration video}}
\par\scriptsize External video: PDF viewer support varies. Keep the video available for separate playback.
\column{0.43\textwidth}
\evidence{validation_demo.png}{Insert terminal comparison or annotated crop screenshot}
\end{columns}
\end{frame}

\begin{frame}{Performance: what the measurements support}
\begin{center}
\begin{tabular}{lr}
\toprule
Measure & Observed result \\
\midrule
Mean of per-crop CPU medians & $52.66\ \mu$s \\
Estimated active FPGA time & $83.84$--$83.92\ \mu$s \\
Mean FPGA round-trip time & $1.250$ s \\
Total components, CPU / FPGA & 24 / 21 \\
Exact agreement & 17 / 20 crops \\
\bottomrule
\end{tabular}
\end{center}
\begin{itemize}
\item FPGA acceleration is \textbf{not demonstrated} by this experiment.
\item JTAG transport dominates wall-clock latency; report it separately from active cycles.
\item Small synthetic crops and failed timing limit performance conclusions.
\item Five passing crops contain no components; more challenging cases are needed.
\end{itemize}
\end{frame}

\begin{frame}{Potential FPGA advantages and current limitations}
\begin{columns}[T]
\column{0.48\textwidth}
\begin{block}{Architectural opportunities}
\begin{itemize}
\item Streaming processing near the map producer.
\item Custom precision and memory organization.
\item Overlap input, processing and output.
\item Avoid unnecessary threshold accesses through margin gating.
\end{itemize}
\end{block}
\column{0.48\textwidth}
\begin{alertblock}{Not yet demonstrated here}
\begin{itemize}
\item Higher throughput or lower latency.
\item Lower power or energy per image.
\item Full-frame scalability.
\item Reliable exact agreement across inputs.
\end{itemize}
\end{alertblock}
\end{columns}
\medskip FPGA benefits depend on architecture, implementation and measurement; using an FPGA alone does not guarantee them.
\end{frame}

\begin{frame}{Next steps: correctness first, then scale and benchmark}
\begin{enumerate}
\item \textbf{Close timing:} inspect critical paths; pipeline or sequence union-find reads; use synchronous block RAM where feasible.
\item \textbf{Check the clock:} a slower experiment needs an actual divided/PLL clock and matching constraints. Editing only the SDC does not slow hardware.
\item \textbf{Reproduce failures:} simulate exact failing maps, merges, border cases, label reuse and backpressure; rerun the board tests.
\item \textbf{Scale coverage:} independent tiles cover a photo but require boundary stitching for full-photo component equivalence.
\item \textbf{Measure fairly:} optimized CPU baseline, verified core clock, repeated board runs, resource use and measured power.
\item \textbf{Use model outputs:} repeat on trained IIM maps and evaluate head localization against annotations separately.
\end{enumerate}
\end{frame}

\begin{frame}{Final assessment}
\begin{block}{Completed}
CPU reference agreement; quantization workflow; synthesizable integration; board programming; working USB-Blaster transfer; reproducible CPU-versus-board validator.
\end{block}
\begin{alertblock}{Current decision}
\textbf{Partial functional success, not a validated accelerator.} The 20-crop run has 3 failures and the design violates its 50 MHz setup requirement. No speedup, energy saving or full-frame correctness claim is justified yet.
\end{alertblock}
\begin{block}{Immediate milestone}
A timing-clean build that passes all existing component comparisons before increasing image size or claiming performance advantages.
\end{block}
\end{frame}

\begin{frame}{Evidence and reproducibility}
\small
\begin{itemize}
\item CPU checks: \texttt{reference\_algorithm/validate.py}.
\item Board comparisons: \texttt{reference\_algorithm/validate\_fpga.py}.
\item Experiment data: \texttt{exp/fpga\_samples\_20.json}.
\item Timing evidence: \texttt{quartus/output\_files/de2\_115\_jtag.sta.summary}.
\item Setup and commands: \texttt{docs/USB\_BLASTER.md}.
\item Board results and table are from the recorded 20-crop run; CPU full-map counts were rechecked during presentation preparation.
\end{itemize}
\end{frame}
\end{document}
